\input{../tex-inputs/latex-leseansicht-vorspann}

\section[Arthur Schnitzler, Richard Beer-Hofmann, Alfred Kerr und Paul Goldmann an Gustav Schwarzkopf, {[}19. 8. 1900{]}]{L04523 Arthur Schnitzler, Richard Beer-Hofmann, Alfred Kerr und Paul Goldmann an Gustav Schwarzkopf, {[}19. 8. 1900{]}}
\nopagebreak\mylabel{L04523v}
\rehead{ }\normalsize\beginnumbering\briefempfaengerindex{Schwarzkopf, Gustav@\textsc{Schwarzkopf, Gustav}!zzzGoldmann, Paul@\emph{von Paul Goldmann}!1900-08-191@{{[}19. 8. 1900{]}}|(be}\briefempfaengerindex{Schwarzkopf, Gustav@\textsc{Schwarzkopf, Gustav}!zzzKerr, Alfred@\emph{von Alfred Kerr}!1900-08-191@{{[}19. 8. 1900{]}}|(be}\briefempfaengerindex{Schwarzkopf, Gustav@\textsc{Schwarzkopf, Gustav}!zzzBeer-Hofmann, Richard@\emph{von Richard Beer-Hofmann}!1900-08-191@{{[}19. 8. 1900{]}}|(be}\briefempfaengerindex{Schwarzkopf, Gustav@\textsc{Schwarzkopf, Gustav}!zzzSchnitzler, Arthur@\emph{von Arthur Schnitzler}!1900-08-191@{{[}19. 8. 1900{]}}|(be}
\toendnotes[C]{\smallbreak\pagebreak[2]}
\correspDesc{Versand  durch Arthur Schnitzler, Richard"> Beer-Hofmann, Alfred Kerr, Paul Goldmann am [19. 8. 1900] in Thusis
\newline{}Übermittlung  am 20. 8. 1900 in Thusis
\newline{}Erhalt  durch Gustav Schwarzkopf im Zeitraum [21. 8. 1900 – 25. 8. 1900?] in Wien}\toendnotes[C]{\smallbreak}
\Standort{DLA, A:Schnitzler, HS.1985.1.1897.}
\physDesc{Bildpostkarte, 248 Zeichen
\newline{}Handschrift Arthur Schnitzler: Bleistift, deutsche Kurrent
\newline{}Handschrift Richard Beer-Hofmann: Bleistift, lateinische Kurrent
\newline{}Handschrift Alfred Kerr: Bleistift, lateinische Kurrent
\newline{}Handschrift Paul Goldmann: Bleistift, deutsche Kurrent
\newline{}Versand: 1) Stempel: »\nobreak{}\oindex{Thusis@\textbf{Thusis}|pwk}Thusis, 20. VIII. 00, VII\nobreak{}«.   2) Stempel: »\nobreak{}\oindex{Wien@\textbf{Wien}|pwk}Wien, \textcolor{gray}{21. 08} 00\nobreak{}«. }\toendnotes[C]{\smallbreak}
\pstart
           \noindent{}\centering{}{\pb}\textcolor{gray}{\textbf{Gruss aus Thusis\oindex{Thusis@\textbf{Thusis}|pw}.}}\pend
           \pstart{}{\pb}Hn \textsc{Gustav Schwarzkopf}\pend{}\pstart{}Wien\oindex{Wien@\textbf{Wien}|pw}\pend{}\pstart{}\textsc{I. Tiefer Graben 23}\oindex{Wien@\textbf{Wien}!I., Innere Stadt@\textbf{I., Innere Stadt}!Tiefer Graben 23@\textbf{Tiefer Graben 23}|pw}.\pend{}{\bigskip}\vspace{1em}
\pstart
           \noindent{}{\pb}\label{K_L04523-1v}\edtext{Heute}{\lemma{\textnormal{\emph{Heute}}}\Cendnote{\textnormal{Der Poststempel datiert den Versand auf den 20. 8. 1900. Aus dem \emph{Tagebuch}\pwindex{Schnitzler, Arthur 15. 5. 1862 Wien – 21. 10. 1931 ebd.@\textsc{Schnitzler, Arthur} (15. 5. 1862 Wien – 21. 10. 1931 ebd.)!Tagebuch@\strich\emph{Tagebuch}|pwk} geht hervor, dass das Schlappiner Joch\oindex{Schlappiner Joch@\textbf{Schlappiner Joch}|pwk} bereits am 19. 8. 1900 überquert wurde, also dürfte die Karte an diesem Tag geschrieben worden sein.}}}\label{K_L04523-1} ſind wir um 4 \textcolor{gray}{a}ufgeſtanden und 
               haben das \textsc{Schlappiner} Joch\oindex{Schlappiner Joch@\textbf{Schlappiner Joch}|pw} überſetzt,
      um in die Schweiz\oindex{Schweiz@\textbf{Schweiz}|pw} zu gelangen.\pend
           
\pstart
           Morgen fahren wir
      nach Pontreſina\oindex{Pontresina@\textbf{Pontresina}|pw}.\pend
           
\pstart
           Herzlichſt Ihr{\\[\baselineskip]}\spacefill\mbox{ArthSchn}\pend
           \leftskip=0em{}\selectlanguage{ngerman}\vspace{1em}
\pstart
           \noindent{}{[}hs. Beer-Hofmann:{]} Herzl. Ihr \spacefill\mbox{Richard}\pend
           \selectlanguage{ngerman}\vspace{1em}
\pstart
           \noindent{}{[}hs. Kerr:{]} \spacefill\mbox{Alfr. Kerr}\pend
           \selectlanguage{ngerman}\vspace{1em}
\pstart
           \noindent{}{[}hs. Goldmann:{]} Herzl. Gruß
               \spacefill\mbox{Dr. Goldmann}\pend
           \selectlanguage{ngerman}\endnumbering\briefempfaengerindex{Schwarzkopf, Gustav@\textsc{Schwarzkopf, Gustav}!zzzGoldmann, Paul@\emph{von Paul Goldmann}!1900-08-191@{{[}19. 8. 1900{]}}|)be}\briefempfaengerindex{Schwarzkopf, Gustav@\textsc{Schwarzkopf, Gustav}!zzzKerr, Alfred@\emph{von Alfred Kerr}!1900-08-191@{{[}19. 8. 1900{]}}|)be}\briefempfaengerindex{Schwarzkopf, Gustav@\textsc{Schwarzkopf, Gustav}!zzzBeer-Hofmann, Richard@\emph{von Richard Beer-Hofmann}!1900-08-191@{{[}19. 8. 1900{]}}|)be}\briefempfaengerindex{Schwarzkopf, Gustav@\textsc{Schwarzkopf, Gustav}!zzzSchnitzler, Arthur@\emph{von Arthur Schnitzler}!1900-08-191@{{[}19. 8. 1900{]}}|)be}\mylabel{L04523h}
\begin{anhang}
\end{anhang}\newcommand{\dateiname}{L04523}\newcommand{\titel}{Arthur Schnitzler, Richard Beer-Hofmann, Alfred Kerr und Paul Goldmann an Gustav Schwarzkopf, [19. 8. 1900]}\newcommand{\editorInnen}{Herausgegeben von Jahnke, SelmaMüller, Martin Anton}\input{../tex-inputs/latex-leseansicht-abspann}
